% ---------------------------------------------------------------------------
% Layout modelled on Murphy, "Machine Learning: A Probabilistic Perspective":
% a large chapter number with an italic title, section numbers hanging in the
% left margin, bold key terms, numbered equations, numbered theorems and
% algorithms, captions that name the script which drew each figure, footnotes
% for asides, and exercises at the end of each chapter.
% ---------------------------------------------------------------------------
\usepackage[paperwidth=7.5in, paperheight=10in, left=1.45in, right=0.85in,
            top=0.95in, bottom=0.95in, headsep=0.3in]{geometry}
\usepackage[T1]{fontenc}
\usepackage[utf8]{inputenc}
\usepackage[sc]{mathpazo}
\linespread{1.06}
\usepackage[scaled=0.88]{helvet}
\usepackage[varqu,scaled=0.95]{zi4}
\usepackage{microtype}
\usepackage{amsmath,amssymb,bm}
\usepackage{amsthm}
\usepackage{graphicx}
\usepackage{booktabs,array,multirow,tabularx,longtable}
\usepackage{xcolor}
\usepackage{enumitem}
\usepackage{caption}
\usepackage{subcaption}
\usepackage{titlesec}
\usepackage{fancyhdr}
\usepackage{listings}
\usepackage{tikz}
\usetikzlibrary{positioning,arrows.meta,fit,calc,shapes.geometric,shapes.misc,backgrounds,
                decorations.pathreplacing,patterns,matrix}
\usepackage{placeins}
\usepackage[hidelinks]{hyperref}
\usepackage{natbib}
\setcitestyle{authoryear,round}

% ---- palette (dataviz reference palette, light mode; text stays in ink) ----
\definecolor{vizblue}{HTML}{2A78D6}
\definecolor{vizorange}{HTML}{EB6834}
\definecolor{vizaqua}{HTML}{1BAF7A}
\definecolor{vizyellow}{HTML}{EDA100}
\definecolor{vizmagenta}{HTML}{E87BA4}
\definecolor{vizgreen}{HTML}{008300}
\definecolor{vizviolet}{HTML}{4A3AA7}
\definecolor{vizred}{HTML}{E34948}
\definecolor{inkprimary}{HTML}{0B0B0B}
\definecolor{inksecondary}{HTML}{52514E}
\definecolor{inkmuted}{HTML}{898781}
\definecolor{gridline}{HTML}{E1E0D9}
\definecolor{codeback}{HTML}{F6F5F2}

% ---- chapter and section headings (numbers hang in the margin) ----
\titleformat{\chapter}[display]
  {\normalfont\raggedright}
  {\fontsize{64}{64}\selectfont\bfseries\color{inkprimary}\thechapter}
  {6pt}
  {\Huge\itshape}
  [\vspace{2pt}]
\titleformat{name=\chapter,numberless}[display]
  {\normalfont\raggedright}{}{0pt}{\Huge\itshape}
\titlespacing*{\chapter}{0pt}{0.2in}{0.45in}
\titleformat{\section}{\normalfont\Large\bfseries}{\llap{\thesection\hspace{1.1em}}}{0pt}{}
\titleformat{\subsection}{\normalfont\large\bfseries}{\llap{\thesubsection\hspace{1.1em}}}{0pt}{}
\titleformat{\subsubsection}{\normalfont\normalsize\bfseries}{\llap{\thesubsubsection\hspace{1.1em}}}{0pt}{}
\titlespacing*{\section}{0pt}{3.2ex plus 1ex minus .2ex}{1.4ex plus .2ex}
\titlespacing*{\subsection}{0pt}{2.6ex plus 1ex minus .2ex}{1.1ex plus .2ex}
\titlespacing*{\subsubsection}{0pt}{2.2ex plus 1ex minus .2ex}{0.9ex plus .2ex}
\setcounter{secnumdepth}{3}
\setcounter{tocdepth}{1}

% ---- running heads: page number, then "Chapter N. Title" in italics ----
\pagestyle{fancy}
\fancyhf{}
\renewcommand{\headrulewidth}{0pt}
\makeatletter
\renewcommand{\chaptermark}[1]{\markboth{\ifnum\value{chapter}>0 \@chapapp\ \thechapter. \fi #1}{}}
\makeatother
\fancyhead[L]{\small\thepage}
\fancyhead[R]{\small\itshape\nouppercase{\leftmark}}
\fancypagestyle{plain}{\fancyhf{}\fancyfoot[C]{\small\thepage}}

% ---- captions: bold label, small text; each ends with its generating script ----
% Captions end in file paths, so they may use loose lines rather than run into the margin.
\DeclareCaptionFont{sloppy}{\sloppy}
\captionsetup{font={small,sloppy}, labelfont=bf, labelsep=colon, justification=justified,
              singlelinecheck=false, skip=6pt}
\newcommand{\generatedby}[1]{Figure generated by \file{#1}.}
\newcommand{\drawnin}[1]{Diagram drawn in \file{#1}.}

% ---- vocabulary ----
\newcommand{\keyterm}[1]{\textbf{#1}}
% File paths in typewriter type, breakable after "/" and "_" so long paths never overflow.
\ExplSyntaxOn
\tl_new:N \l__sol_path_tl
\NewDocumentCommand \file { m }
  {
    \tl_set:Nn \l__sol_path_tl {#1}
    \tl_replace_all:Nnn \l__sol_path_tl { / } { / \allowbreak }
    \tl_replace_all:Nnn \l__sol_path_tl { \_ } { \_ \allowbreak }
    \texttt { \upshape \l__sol_path_tl }
  }
\ExplSyntaxOff
% A third line-breaking pass with extra stretch before any line is allowed to overflow.
\setlength{\emergencystretch}{3em}
\newcommand{\Ha}{H$\alpha$}
\newcommand{\PQ}{\mathrm{PQ}}
\newcommand{\SQ}{\mathrm{SQ}}
\newcommand{\RQ}{\mathrm{RQ}}
\newcommand{\IoU}{\mathrm{IoU}}
\newcommand{\TP}{\mathrm{TP}}
\newcommand{\FP}{\mathrm{FP}}
\newcommand{\FN}{\mathrm{FN}}
\newcommand{\R}{\mathbb{R}}
\newcommand{\E}{\mathbb{E}}
\newcommand{\argmin}{\operatorname*{arg\,min}}
\newcommand{\argmax}{\operatorname*{arg\,max}}
\newcommand{\sigmoid}{\sigma}
\newcommand{\softmax}{\operatorname{softmax}}

% ---- code listings ----
\lstdefinestyle{repo}{
  language=Python,
  basicstyle=\ttfamily\footnotesize,
  keywordstyle=\color{vizblue}\bfseries,
  commentstyle=\color{inkmuted},
  stringstyle=\color{vizorange!85!black},
  numberstyle=\tiny\color{inkmuted},
  numbers=left, numbersep=10pt,
  backgroundcolor=\color{codeback},
  frame=l, framerule=1.2pt, rulecolor=\color{vizblue!45},
  xleftmargin=20pt, framexleftmargin=4pt,
  showstringspaces=false, breaklines=true, breakatwhitespace=true,
  columns=fullflexible, keepspaces=true, tabsize=4,
  aboveskip=8pt, belowskip=6pt,
  % Listings are kept to ASCII: pdflatex's listings cannot typeset multi-byte UTF-8.
}
\lstset{style=repo}
\renewcommand{\lstlistingname}{Listing}
\captionsetup[lstlisting]{font=small, labelfont=bf, skip=4pt, justification=raggedright,
                           singlelinecheck=false}

% ---- theorems, as in Murphy: numbered within the chapter ----
\theoremstyle{plain}
\newtheorem{theorem}{Theorem}[chapter]
\newtheorem{proposition}[theorem]{Proposition}
\theoremstyle{definition}
\newtheorem{definition}[theorem]{Definition}

% ---- algorithms, as in Murphy: a ruled block with numbered lines ----
% Usage: \begin{algorithmblock}{title} \item first line \item \Ind indented line \end{algorithmblock}
\newcounter{algorithmblock}[chapter]
\renewcommand{\thealgorithmblock}{\thechapter.\arabic{algorithmblock}}
\def\theHalgorithmblock{\thechapter.\arabic{algorithmblock}}   % unique hyperref anchors
\newcommand{\Ind}{\hspace*{1.4em}}
\newenvironment{algorithmblock}[1]{%
  \par\medskip\refstepcounter{algorithmblock}%
  \noindent\rule{\linewidth}{0.7pt}\par\nopagebreak\vspace{0.3ex}%
  \noindent{\small\textbf{Algorithm \thealgorithmblock:}\ #1}\par\nopagebreak\vspace{-0.4ex}%
  \noindent\rule{\linewidth}{0.4pt}\par\nopagebreak\vspace{-0.6ex}%
  \small
  \begin{enumerate}[label={\color{inkmuted}\scriptsize\arabic*}, leftmargin=2.2em, labelsep=0.9em,
                    itemsep=0.5pt, parsep=0pt, topsep=3pt]}%
  {\end{enumerate}\vspace{-0.6ex}\noindent\rule{\linewidth}{0.7pt}\par\medskip}

% ---- epigraph under a chapter's first section, as in Murphy ----
\newcommand{\epigraph}[2]{\begin{quote}\small\itshape #1\par\nopagebreak\vspace{0.3ex}
  {\raggedleft\upshape--- #2\par}\end{quote}\vspace{0.5ex}}

% ---- exercises, as in Murphy: "Exercise N.M Title", optional source line ----
\newcounter{exercise}[chapter]
\renewcommand{\theexercise}{\thechapter.\arabic{exercise}}
\def\theHexercise{\thechapter.\arabic{exercise}}   % unique hyperref anchors
\newenvironment{exercise}[2][]{%
  \par\medskip\refstepcounter{exercise}\noindent\textbf{Exercise \theexercise}\ \ #2%
  \ifx\relax#1\relax\else\par\noindent{\small(Source: #1.)}\fi\par\smallskip\noindent\ignorespaces}%
  {\par\medskip}

% ---- TikZ styles shared by the diagrams ----
\tikzset{
  >={Stealth[length=2.2mm]},
  box/.style={draw=#1, line width=0.8pt, fill=#1!9, rounded corners=2.5pt, align=center,
              inner sep=4pt, font=\sffamily\footnotesize, text=inksecondary},
  box/.default=inkmuted,
  flow/.style={->, line width=0.8pt, draw=inksecondary},
  dashflow/.style={->, line width=0.7pt, draw=inkmuted, dashed},
  note/.style={font=\sffamily\scriptsize, text=inksecondary, align=center},
  title/.style={font=\sffamily\footnotesize\bfseries, text=inkprimary},
}

\hypersetup{colorlinks=true, linkcolor=vizblue!80!black, citecolor=vizblue!80!black,
            urlcolor=vizblue!80!black}
